% PUBLIC PROVENANCE SNAPSHOT — TOPIC AND EQUATION MAP ONLY
%
% Machine-local upstream files:
% /Users/setup/Library/CloudStorage/Box-Box/Teaching/6050/LaTeX/Module 12/12.1S Multimodal.tex
% SHA-256: 21c0751a49b44efe1c1c21e38246ca1b68575bb100c05704cfe2d32de32caaf4
% /Users/setup/Library/CloudStorage/Box-Box/Teaching/6050/LaTeX/Module 12/12.0 Multimodal_GenAI.tex
% SHA-256: 52e0e1f1239d44262838653419f753b612bf1538ccbe5003d43e1fb9cda13dc8
%
% The Module 12 course archive contains no matching recorded-video transcript.
% The dedicated 12.1 deck is therefore the chapter's instructor-voice spine; the
% longer 12.0 article supplies only its matching multimodality and evaluation map.
%
% Omitted from this public snapshot: the local Beamer preamble; all diagrams,
% external images, product examples, executable code, Colab links, current-model
% surveys, unverified benchmark numbers, scale/performance prescriptions, and
% categorical robustness, generalization, cognition, or social-impact claims.
% Those omissions include several statements that would require more qualification
% than a provenance snapshot can supply. No D2L material has been retained.
%
% Book prose, derivations, code, experiments, and figures are independently written
% and checked against the primary sources named in Chapter 20.

\documentclass{article}
\usepackage{amsmath,amssymb}
\begin{document}

\section{Instructor-authored multimodal spine retained for traceability}

\begin{enumerate}
  \item Vision and language have different raw representations. A multimodal model
  needs a declared relation between them rather than an assumption that their
  coordinates are already comparable.
  \item Two encoders map paired observations to a shared embedding dimension. The
  pair relation, not a class label, supplies the training correspondence.
  \item In a batch of $N$ pairs, every image embedding is compared with every text
  embedding, producing an $N\times N$ score matrix. Diagonal entries are declared
  matches; off-diagonal entries act as within-batch alternatives.
  \item Cosine similarity compares directions after norm control:
  \[
    \operatorname{sim}(\mathbf z_i,\mathbf z_t)
    =\frac{\mathbf z_i^\top\mathbf z_t}
    {\lVert\mathbf z_i\rVert\lVert\mathbf z_t\rVert}.
  \]
  \item The image-to-text contrastive objective is a softmax classification loss
  over each score row:
  \[
    \mathcal L_{i\rightarrow t}
    =-\frac1N\sum_{i=1}^{N}
    \log\frac{\exp(s_{ii}/\tau)}
    {\sum_{j=1}^{N}\exp(s_{ij}/\tau)}.
  \]
  The reverse direction treats each score column as the candidate set, and the
  symmetric objective averages the two directions:
  \[
    \mathcal L_{\mathrm{sym}}
    =\tfrac12(\mathcal L_{i\rightarrow t}+\mathcal L_{t\rightarrow i}).
  \]
  \item Cross-modal retrieval ranks candidates in one modality using a query from
  the other. Recall at $k$ checks whether the declared mate occurs among the first
  $k$ candidates.
  \item Zero-shot classification is a bounded retrieval construction: encode a
  fixed candidate set of text descriptions and select the description with largest
  image--text similarity. It does not by itself generate text or prove that an
  arbitrary unseen concept was represented during training.
  \item Contrastive matching, binary image--text matching, and autoregressive
  captioning are different objectives. A dual encoder that ranks existing candidates
  is not thereby a generative vision--language model.
  \item Evaluation must examine the pair definition, candidate pool, duplicate or
  ambiguous matches, distribution shifts, subgroup behavior, compositional failures,
  privacy, consent, and data provenance rather than reporting one aggregate score.
\end{enumerate}

\section{Primary references retained from the lecture map}

\begin{itemize}
  \item Radford et al. (2021), ``Learning Transferable Visual Models From Natural
  Language Supervision.''
  \item Jia et al. (2021), ``Scaling Up Visual and Vision-Language Representation
  Learning With Noisy Text Supervision.''
  \item van den Oord, Li, and Vinyals (2018), ``Representation Learning with
  Contrastive Predictive Coding.''
\end{itemize}

\end{document}
